\scanpage{7}
D'après les sorites généraux sur les foncteurs représentables, $P\otimes_A M$
dépend fonctoriellement de $\handwritten{(}M\handwritten{,P)}$ \deleted{et de $P$,} \deleted{là où il} pour les \inserted{\handwritten{\unsure{couples|couple} d'}}arguments pour
lesquels il est défini. De plus, on voit tout de suite que si on a une
suite exacte
\[
M''\handwritten{\to}M'\handwritten{\unsure{-|\to}}M\handwritten{\unsure{\to|-}}0
\]
de $A$-modules à gauche, telle que $P\otimes_A M'$ et $P\otimes_A M''$ existent, alors $P\otimes_A M$
existe et \deleted{on a une}\inserted{\handwritten{la}} suite \deleted{exacte}
\[
P\otimes_A M''\handwritten{\to}P\otimes_A M'\handwritten{\unsure{-|\to}}P\otimes_A M\handwritten{\unsure{-|\to}}0
\]
\inserted{\handwritten{est exacte}} (en particulier, $P\otimes_A M$ est \uline{foncteur exact à droite en} $M$, sous réserve
d'existence).
% The block beginning "On a le fait analogue" is enclosed and carried down by an arrow to "Combinant".
\deleted{Cela prouver} Plus généralement, le foncteur
$P\otimes_A M$ en $M$, ou en $P$, commute aux limites inductives quelconques (pas nécessairement filt\unsure{r}antes), au sesn suivant: si les $P\otimes_A M_i$ existent, \deleted{et $\lim$}
alors $P\otimes_A\mathop{\underset{\handwritten{\longrightarrow}}{\lim}}\limits_{\handwritten{i}}M_i$ existe sss $\underset{\handwritten{\longrightarrow}}{\lim}P\otimes_A M_i$ existe, et les deux sont can isomorphes. Ceci s'applique en particulier aux limites inductives sur un ensemble discret, i.e. aux sommes directes.
\inserted{On a le fait analogue pour des suites exactes en $P$\inserted{\handwritten{, et plus gén. des $\varinjlim$ en $P$}}. D'autre
part, $P\otimes_A A_g$ existe toujours et\handwritten{\unsure{/}}on a un isomorphisme canonique
\[
P\otimes_A A_g\handwritten{\simeq}P\;.
\]
\deleted{(Pour}(Tous ces faits sont triviaux sur les définitions, en interprétant
en termes de foncteurs en $Q$).}
Combinant ces remarques, on voit
que $P\otimes_A M$ existe \deleted{toujours} dès que $M$ est de présentation finie (car il
existe pour $M=A_g^{(I)}$, $I$ ensemble fini, donc pour des conoyaux d'homomorphismes $A_g^{(I)}\handwritten{\longrightarrow}A_g^{(J)}$, $I$ et $J$ finis), et sans restriction sur $M$ si dans $\underline C'$
les sommes directes quelconques existent, (même \deleted{\unsure{ar}}argument, sans restriction
de finitude sur $I,J$). De plus, on a un pro\deleted{\unsure{t}}cédé de construction de $P\otimes_A M$
par ``genérateurs et relations'' en termes de ``générateurs et relations''
pour $M$, i.e. en termes \deleted{d'un} de la matrice définissant une présentation
\deleted{finie} de $M$ (cette matrice définit une matrice sur $I\mathbin{\mathrm{x}}J$\handwritten{,} définissant une
application $P^{(I)}\handwritten{\longrightarrow}P^{(J)}$, dont le conoyau est le produit tensoriel cherché).
Tout ceci devrait figurer dans un numéro à part\inserted{ au Chap $0$}, sans trainer l'accent '
sur la catégorie abélienne. Du point de vue notations, si $C'$ est \deleted{la duale}\inserted{l'opposée}
d'une catégorie abélienne $C_1$, \deleted{on note $\operatorname{Hom}$} les $A$-$C'$-modules à droite
sont les $A$-$C_1$-modules à gauche, et le produit tensoriel dans $C'$ se note
(en termes de $C_1$)
\[
\operatorname{Hom}_A(M,P)
\]
défini directement par
\[
\deleted{\operatorname{Hom}_{Q}(\operatorname{Hom}(M,P))}
\]
% archive sheet number: 6

